% Shortened Chapter 3, part 10 of 10: section 3.10.
% Include after section 3.9; do not include alongside the original section 3.10.
% All six subsection slots, all seven labels and all TWELVE rows of the
% divergence ledger are retained. Chapters 2, 4, 5 and 6 reference
% sec:synthesis-and-gap, sec:four-limitations, sec:research-gap,
% sec:addressing-gap and sec:divergence-ledger.

\section{Synthesis and Research Gap}
\label{sec:synthesis-and-gap}

This section consolidates the gap statements closing \autoref{sec:corpus}--\autoref{sec:imbalance-review}: what the corpus collectively establishes, the four limitations recurring across every topic, the resulting research gap, and which parts of it this thesis closes. It also collects in one place every declared divergence between the system built here and the literature it draws on.

\subsection{What the corpus establishes}
\label{sec:corpus-establishes}

Read as a whole, the reviewed literature settles a good deal. \textbf{On data:} the corpora are small, demographically narrow, laboratory-bound and irreducibly imbalanced by what can be elicited, and CASME II remains the primary single-corpus benchmark under its native leave-one-subject-out protocol (\autoref{sec:corpus}, \autoref{sec:imbalance-review}). \textbf{On representation:} motion rather than appearance is the appropriate input; a strain field adds a further cue, though on CASME II only weakly and under a protocol that is not subject-disjoint; and both the magnification factor and the interpolated sequence length have interior optima located by sweep (\autoref{sec:evm-review}--\autoref{sec:tim-review}). \textbf{On architecture:} at this data size capability bought without parameters outperforms capability bought with them, and self-attention is the strongest-performing mechanism reported in the corpus (\autoref{sec:cnn-review}--\autoref{sec:transformer-review}). \textbf{On evaluation:} accuracy alone is inadequate under skew and must give way to macro-averaged measures accumulated across folds before averaging (\autoref{sec:megc-metrics}, \autoref{sec:imbalance-metrics}).

\subsection{Four limitations that recur across every topic}
\label{sec:four-limitations}

The per-section gap statements are not nine independent observations. Four structural problems generate most of them.

\textbf{(a) Components are isolated one at a time, inside one proposed design, and never jointly.} Matched-pair measurement is not absent from the corpus. On the input side \autoref{tab:evm-effect} isolates magnification and \autoref{tab:strain-performance} isolates strain; on the architectural side Xia et al. compare a basic recurrent convolutional network against three variants adding a wide expansion, a shortcut connection and an attention unit, holding every other parameter constant across all four (\autoref{sec:cnn-limited-data}), and separately sweep depth against input resolution (\autoref{sec:cnn-shrink}); SLSTT's Mean-versus-LSTM comparison isolates its aggregator (\autoref{sec:slstt-results}).

What no reviewed study does is vary \textit{several} components together, so that their interactions can be read. Every ablation in the corpus is internal to one proposed design and moves one factor at a time from that design's own defaults; across papers, systems are compared whole --- STSTNet against OFF-ApexNet, SLSTT against STSTNet --- so a difference in the reported number cannot be attributed to any one stage. There is therefore no basis for deciding which of the learned stages are load-bearing, nor whether one component's value depends on another's presence.

\textbf{(b) The evaluation protocol is unstable, and under-reported.} Sample counts, label-set size and the choice among leave-one-subject-out, leave-one-video-out and leave-one-sample-out all vary inconsistently, and the field's most-quoted baseline is reported under two different protocols by its two primary sources and under two more within the originating paper itself (\autoref{sec:original-baseline}). Given this corpus's fold composition (\autoref{sec:corpus-threats}), these are not bookkeeping differences: they change what the numbers mean.

\textbf{(c) The strongest results depend on regimes a single-corpus study cannot access.} The best CASME II figures in the corpus come from composite-database training (\autoref{sec:megc-metrics}) and, in SLSTT's case, ImageNet pre-training (\autoref{sec:transformer-limitations}) --- neither a property of the architecture being credited. Almost no evidence exists about what these components do when trained from scratch on one small corpus, which is the situation most practitioners are in.

\textbf{(d) Preprocessing parameters have interior optima, and are transferred without re-validation.} The magnification factor, the interpolated sequence length and the input resolution all rise then fall, all have their optima located only by sweep, and none has a principled selection rule (\autoref{sec:evm-evidence}, \autoref{sec:tim-how-many}, \autoref{sec:cnn-shrink}). Values are carried between pipelines whose downstream models differ entirely: $\alpha = 30$ moves from a single-frame appearance encoder to other settings, TIM10 propagates from one baseline paper to the whole field, and $\lambda = 10^{-4}$ transfers from CIFAR though its own authors re-tuned it for ImageNet (\autoref{sec:simam-evidence}).

\subsection{The research gap}
\label{sec:research-gap}

Taking (a)--(d) together:

\textbf{The field has established a standard pipeline for micro-expression recognition --- magnify, compute motion, normalise duration, extract spatially, model temporally --- but has never measured which of its stages are responsible for the result. Because components are only ever varied one at a time inside a single proposed system, under protocols that vary in unstated ways, and predominantly in training regimes richer than a single small corpus provides, it is not currently possible to say which parts of that pipeline earn their place, which are inherited convention, and which depend on the presence of another.}

This matters practically as well as scientifically. A practitioner assembling this pipeline today has no published basis for allocating effort: for judging whether the most computationally expensive stage is the most valuable, whether a component with a negligible average effect may still matter in a specific failure, or whether a ranking read off an incomplete evaluation of a corpus this size is stable at all.

\subsection{How this thesis addresses the gap --- and what it does not}
\label{sec:addressing-gap}

\textbf{What it does.} This thesis treats the pipeline as a factorial experiment rather than a proposal: four components varied independently across every architecturally valid cell of a 2\textsuperscript{4} matrix, twelve configurations, each evaluated under \textbf{complete 25-fold leave-one-subject-out on CASME II} with every non-varied factor held identical. That gives the joint, multi-component measurement (a) says the corpus lacks; reporting the protocol in full addresses (b) for this study at least; and training from scratch is the regime (c) identifies as unevidenced. Each component's marginal contribution is read from pairs of runs differing in one factor alone, and the distribution of those four effects --- how much of the pipeline's performance each stage accounts for, whether any subtracts, and what each costs --- is the substantive answer to the gap.

Three further contributions follow from the review rather than from the experiment:

\begin{itemize}
	\item \textbf{A caution about partial evaluations.} This project's earlier evaluations, on smaller held-out sets and incomplete fold counts, ranked the twelve configurations differently from the complete run reported here. They are not evidence of protocol sensitivity, since they differ in sample count, training budget and class coverage too, and the earliest carry a data-routing defect --- but the episode illustrates limitation (b) in practice, and only the complete run is reported.
	\item \textbf{A repaired measurement.} That defect left the magnification switch inert, producing bit-identical results across every pair those runs contained. What made it detectable as a defect rather than a null result was the review's evidence that magnification has a clear, repeatable effect (\autoref{sec:evm-implications}).
	\item \textbf{Two opposite class-collapse failures.} A single loss-level correction suppressed the minority class; adding a data-level correction on top of it suppressed the majority class instead. No reviewed paper reports either collapse (\autoref{sec:imbalance-double-correction}).
\end{itemize}

\textbf{What it does not do.} Three of the four limitations are only partially addressed. On (b), this study reports its own protocol completely but cannot repair the field's inconsistency, and its own N = 156 differs from MEGC's 145. On (c), it establishes what these components do from scratch on one corpus, which is a different question from --- not an answer to --- how they behave under composite training; no cross-dataset claim is made anywhere. On (d), the thesis inherits parameter values rather than sweeping them: $\alpha = 10$, T = 32, $\lambda = 10^{-4}$ and $\gamma = 2.0$ are fixed throughout, so the ablation measures architecture at one point in preprocessing space, not across it.

\subsection{A consolidated ledger of declared divergences}
\label{sec:divergence-ledger}

Reviewing the literature closely enough to write this chapter surfaced twelve points at which the implemented system departs from the work it draws on. Several were discovered only through the review. They are collected in \autoref{tab:divergence-ledger} so that no result rests on an unstated assumption.

\begin{table}[htbp]
	\centering
	\footnotesize
	\begin{tabular}{@{}p{0.85cm} p{6.0cm} p{6.6cm}@{}}
		\hline
		\textbf{\S} & \textbf{Divergence} & \textbf{Consequence} \\
		\hline
		3.1.3 & Working set is N = 156, including sadness and fear in Negative; MEGC's CASME II subset is N = 145 & Comparisons against challenge figures are not exactly like-for-like \\
		3.1.2 & The corpus's own registered release is not used; the pipeline reads the original recorded frames, unregistered and uncropped & No landmarking or alignment of any kind precedes the network; tolerance to pose and scale must be learned from 156 clips \\
		3.2.6 & Temporal band is the unnarrowed 5--25 Hz duration rule, where Bai et al. narrowed it to 15--25 Hz as noisy & The pass-band admits the low-frequency range they rejected \\
		3.4.3 & Strain magnitude counts shear once, not twice as in the published Frobenius norm & Shear weighted lower by $\sqrt{2}$ relative to normal strain; untested \\
		3.4.7 & No dedicated strain filtering, where the literature applies Wiener and Gaussian filters & Spatial differentiation can amplify flow-estimation noise, and nothing downstream suppresses it \\
		3.5.7 & Temporal normalisation is uniform index sampling, \textbf{not} manifold-based TIM; no frames are synthesised & Avoids fabricating motion; forfeits the ability to lengthen short clips \\
		3.5.7 & Magnification is applied \textbf{after} subsampling, at a nominal 200 fps & Realised pass-band is clip-dependent and below the intended 5--25 Hz \\
		3.6.6 & Backbone kernels are (1, 3, 3) --- \textbf{no temporal mixing} & The ablation tests a learned \textit{spatial} stem, not spatio-temporal convolution \\
		3.6.6 & Backbone input is 224 $\times$ 224, against Xia et al.'s guidance of $\leq$ 100 $\times$ 100 --- though Bai et al. and Li, Huang and Zhao both use 224 $\times$ 224 & Whatever the stem is measured to contribute may be a resolution effect as much as an architecture effect \\
		3.7.7 & SimAM statistics pooled over the whole clip rather than per frame as in the source; $\lambda$ not re-tuned & Neurons judged against the clip, not the frame; defensible but untested \\
		3.8.6 & Six divergences from published SLSTT (temporal not spatial attention, mean not LSTM aggregation, from scratch, smaller, sinusoidal encoding, short-term flow) & Variable D is a temporal transformer encoder, \textbf{not} a reproduction of SLSTT \\
		3.9.7 & Label smoothing of 0.05, named in the corpus only by Dosovitskiy et al. and never examined on this task & A choice the micro-expression corpus supports neither way, held constant and unmeasured \\
		\hline
	\end{tabular}
	\caption[Every declared departure from the reviewed literature]{Every declared departure from the reviewed literature, with its consequence.}
	\label{tab:divergence-ledger}
\end{table}

Three of these --- the strain formula, the magnification ordering, and the input resolution --- were identified by the review itself rather than being known design decisions, and each is declared as proposed further work. None invalidates the ablation, since every configuration shares the same choices. What they constrain is the \textit{wording} of the conclusions: what is measured is the contribution of a spatial stem at 224 $\times$ 224 and of temporal self-attention in this pipeline, not a verdict on 3D convolution or on SLSTT.

\subsection{What this chapter contributes to the thesis}
\label{sec:ch3-contribution}

Three things carry forward. First, the \textbf{methodology that follows is determined by this review rather than chosen freely}: the corpus dictates the class grouping, representation, fold structure and metric, and each component's parameters come from published sweeps (\autoref{sec:corpus-implications}). Second, the \textbf{interpretation of the central result is prepared here rather than improvised later}: flow and strain already perform spatio-temporal feature extraction analytically (\autoref{sec:flow-implications}, \autoref{sec:strain-implications}), so a negative contribution from a spatial-only backbone at an oversized input resolution is expected rather than anomalous (\autoref{sec:cnn-implications}). Third, the \textbf{limits of what can be claimed are fixed before any number is reported}: a small, single-seed, demographically narrow corpus measured against richer published regimes, with the twelve divergences of \autoref{tab:divergence-ledger} standing throughout --- consequences of this review, not concessions made after the fact.
